\scanpage{17}

\mspage{\handwritten{2}}

(viewed as a component of such a ``generalised boundary''). We now should \deleted{should} get a modular space for ``compact conformal oriented surfaces with generalized boundary'' (of type $g,\handwritten{\nu}$ and rigidified via $X_0$), call it $\overset{\handwritten{\sim}}{\mathrm{MB}}_{\handwritten{g,\nu}}$, homeomorphic to $\underline{\mathrm C}^d\mathbin{\mathrm x}(\underline{\mathrm R}^{\deleted{\unsure{*}}+})^{\handwritten{\nu}}$, where now the second factor corresponds to \deleted{\unsure{th}} the ``parameters'' $\handwritten{\ell_i}$, which are allowed to take also value 0 (which means that the corresponding component of the generalized boundary is just one point). Thus $\overset{\handwritten{\sim}}{\mathrm{MB}}_{\handwritten{g,\nu}}$ appears as a variety with boundary (in the topological sense - in the real analytic sense, the ``boundary'' admits ``corner-like'' points obviously), and $\overset{\handwritten{\sim}}{\mathrm M}_{g,\handwritten{\nu}}$ appears as a part of the boundary.

My interest is in a better geometric understanding of the situation, which should be ``intrinsic'' namely not depend on any particular choice of a surgical decomposition of the reference surface $X_0$ into ``pants'', used in order to describe in a handy way \deleted{the} standard ``coordinate\deleted{\unsure{s}} functions'' on the modular space $\overset{\handwritten{\sim}}{\mathrm{MB}}_{g,\handwritten{\nu}}$. There appears to be a geometrically meaningful retraction of $\overset{\handwritten{\sim}}{\mathrm{MB}}_{g,\handwritten{\nu}}$ \deleted{$\unsure{\mathrm M}$} upon $\overset{\handwritten{\sim}}{\mathrm M}_{g,\handwritten{\nu}}$ (commuting to the operations of the Teichmüller modular group), the fibers being homeomorphic to $(\underline{\mathrm R}^{+})^{\deleted{\unsure{I}}\handwritten{\nu}}$ - more specifically, I expect the semi-group $(\underline{\mathrm R}^+)^{\mathrm I}$ (where I is the set of indices for the ``holes'' of $X_0$) to act on $\mathrm{MB}$ in a natural way, with free action of the subgroup $(\mathrm R^{+\deleted{\unsure{x|*}}})^{\handwritten{\nu}}$ upon $\overset{\handwritten{\sim}}{\mathrm{MB}}^{\mathrm o}$, in such a way that $\overset{\handwritten{\sim}}{\mathrm M}$ is just the quotient of $\overset{\handwritten{\sim}}{\mathrm{MB}}$ by this action (or of $\overset{\handwritten{\sim}}{\mathrm{MB}}^{\mathrm o}$ by the action of the \deleted{\unsure{acting}} corresponding subgroup), and that \deleted{\unsure{e}}each fiber \inserted{$\mathrm F$} is isomorphic to $(\underline{\mathrm R}^+)^{\mathrm I}$ by the choice of any ``origin'' in $\mathrm F\handwritten{\cap}\overset{\handwritten{\sim}}{\mathrm{MB}}^{\mathrm o}$.

Of course, ``comput\deleted{\unsure{i|a}}\unsure{a|i}nnally'', in terms of a decomposition of $X_0$ into pants, the idea of such an operation is pretty obvious - namely letting the components $\handwritten{\lambda}_{\mathrm i}$ of $\handwritten{\lambda\in}(\underline{\mathrm R}^+)^{\mathrm I}$ act as a ``multiplier'' on the corresp\inserted{ond}ing coordinate $\deleted{\unsure{h}}\handwritten{\ell}_{\mathrm i}$. However, it is not clear that this operation is intrinsic - and if it were intrinsic, an intrinsic geometric description would still be \deleted{\unsure{in|is}} desired.

Of course, in the description of the situation proposed above, the retraction of $\overset{\handwritten{\sim}}{\mathrm{MB}}$ upon $\overset{\handwritten{\sim}}{\mathrm M}$ is obtained by multiplying with the 0 multiplier (all $\handwritten{\lambda}_{\mathrm i}$ are 0). Now there is a direct geometrical description of \deleted{this} a retraction, by hyperbolic surgery. Namely, for any \inserted{compact} conformal surface of type $g,\handwritten{\nu}$ with generalized boundary, let's ``fill in'' the holes which correspond to ordinary components of the boundary, which are riemanian oriented circles, by ``gluing in'' the cones on these circles (which are canonically endowed with a conformal structure, using the riemanian structure on the given circles). Thus we get a ``functor'' from \inserted{compact} conformal surfaces with generalized boundary \deleted{\unsure{t}}(of type $g,\handwritten{\nu}$) to compact conformal surfaces \uline{without boundary}, endowed with a system of \handwritten{$\nu$} points (making up a
% archive sheet number: 16
